\pnum{5.9}{Including spin in low-lying two-particle well levels}
Take single-particle well energies $n^2\epsilon$,
$\epsilon=\pi^2\sh^2/(2ma^2)$, and
$\phi_n=\sqrt{2/a}\sin(n\pi x/a)$.
For fermions of spin $1/2$, antisymmetric
\emph{total} states pair a symmetric
spatial factor $\Phi^+_{nm}$ with the
spin singlet $\ket{0,0}$, and antisymmetric
$\Phi^-_{nm}$ with one of the three triplets
$\ket{1,M}$, $M=-1,0,1$.
When $n=m$, only the symmetric spatial
state survives, and there is just one
spin singlet. The first four distinct
energies and degeneracies are therefore
\begin{equation}
\boxed{
\begin{array}{c|c|c}
E/\epsilon&\text{states}&g\\\hline
2&\phi_1(1)\phi_1(2)\ket{0,0}&1\\
5&\Phi^+_{12}\ket{0,0}\text{ and }
\Phi^-_{12}\ket{1,M}&4\\
8&\phi_2(1)\phi_2(2)\ket{0,0}&1\\
10&\Phi^+_{13}\ket{0,0}\text{ and }
\Phi^-_{13}\ket{1,M}&4
\end{array}}.
\end{equation}
For hypothetical \emph{spin-1 bosons},
the two-spin exchange-symmetric sectors
have total $S=0$ (one state) and $S=2$
(five states), while the antisymmetric
sector has $S=1$ (three states).
A symmetric total bosonic wavefunction
therefore pairs $\Phi^+_{nm}$ with $S=0,2$
and $\Phi^-_{nm}$ with $S=1$.
For the same four energy levels, the
degeneracies are
\begin{equation}
\boxed{g(2\epsilon)=6,\quad
g(5\epsilon)=9,\quad
g(8\epsilon)=6,\quad
g(10\epsilon)=9}.
\end{equation}
These counts include spin and presume no
spin-dependent interaction.

\pnum{5.10}{Why three spin halves have no antisymmetric spin state}
A spin-$1/2$ particle has a two-dimensional
spin space $\mathbb C^2$. An entirely
antisymmetric spin state of $N$ particles
belongs to the exterior power
$\bigwedge^N\mathbb C^2$, of dimension
$\binom{2}{N}$. For $N=3$ that dimension
vanishes: $\boxed{\bigwedge^3\mathbb C^2=\{0\}}$.
Equivalently, every triple of arrows
$\uparrow,\downarrow$ contains two equal
arrows; swapping those must reverse the
sign of a completely antisymmetric state
but leaves the arrow assignment unchanged,
so its coefficient is zero.

For three spin-half fermions in a square
well, one spin orbital can be occupied
only once. The three lowest available
spin orbitals are $1\uparrow$, $1\downarrow$,
and either $2\uparrow$ or $2\downarrow$.
Let $\zeta_{n\sigma}(x,s)=\phi_n(x)\chi_\sigma(s)$.
The normalized ground states are the two
Slater determinants
\begin{equation}
\boxed{\Psi_{\sigma}=
\frac1{\sqrt{3!}}
\det\begin{pmatrix}
\zeta_{1\uparrow}(1)&\zeta_{1\downarrow}(1)&\zeta_{2\sigma}(1)\\
\zeta_{1\uparrow}(2)&\zeta_{1\downarrow}(2)&\zeta_{2\sigma}(2)\\
\zeta_{1\uparrow}(3)&\zeta_{1\downarrow}(3)&\zeta_{2\sigma}(3)
\end{pmatrix}},
\quad \sigma=\uparrow,\downarrow.
\end{equation}
Their energy is
$\boxed{E_{\rm ground}=2\epsilon+\epsilon+4\epsilon
=6\epsilon}$ and spin degeneracy is
$\boxed{2}$. In general the determinant
cannot be factored into one spatial
wavefunction times one spin wavefunction.

\pnum{5.11}{Laughlin wavefunction, Slater determinants, and correlation}
Ignoring an overall normalization, the
Laughlin state of $N$ electrons has the form
\begin{equation}
\Psi_q(z_1,\ldots,z_N)=
\left[\prod_{j<k}(z_j-z_k)^q\right]
\exp\left[-\frac12\sum_j|z_j|^2\right],
\qquad q=1,3,5,\ldots .
\end{equation}
Exchanging $z_j,z_k$ changes one factor's
sign by $(-1)^q=-1$; all other factors
are permuted in pairs. The Gaussian is
symmetric. Thus
$\boxed{P_{jk}\Psi_q=-\Psi_q}$.
For $q=1,N=3$, the Vandermonde formula is
\begin{equation}
\prod_{j<k}(z_j-z_k)=
-\det\begin{pmatrix}
1&z_1&z_1^2\\
1&z_2&z_2^2\\
1&z_3&z_3^2
\end{pmatrix}.
\end{equation}
Up to normalization and a global sign,
this is one Slater determinant formed from
single-particle orbitals
$\boxed{\varphi_m(z)\propto z^m\ee^{-|z|^2/2},
\quad m=0,1,2}$.
For $q=3,N=2$,
\begin{align}
(z_1-z_2)^3
&=(z_1^3-z_2^3)-3(z_1^2z_2-z_1z_2^2)\\
&=-\det\begin{pmatrix}1&z_1^3\\1&z_2^3\end{pmatrix}
-3\det\begin{pmatrix}z_1&z_1^2\\z_2&z_2^2\end{pmatrix}.
\end{align}
The state is a superposition of
determinants with different orbital
occupancies $(0,3)$ and $(1,2)$.
The Jastrow factor has a third-order
zero when particles meet, encoding
correlation not captured by a single
determinant. The normalized determinant
coefficients require the normalization
of the $z^m$ orbitals.

\pnum{5.12}{Symmetry projection for a general atomic eigenstate}
For $Z$ particles let $U_\pi$ permute their
position arguments. Since the identical-electron
Hamiltonian commutes with every permutation,
$[H,U_\pi]=0$. Define
\begin{equation}
\Sop=\frac1{Z!}\sum_{\pi\in S_Z}U_\pi,
\qquad
\Aop=\frac1{Z!}\sum_{\pi\in S_Z}
\operatorname{sgn}(\pi)\,U_\pi.
\end{equation}
If $H\psi=E\psi$, then
$H(\Sop\psi)=E(\Sop\psi)$ and
$H(\Aop\psi)=E(\Aop\psi)$.
These states must be normalized
if nonzero. If $\psi$ is symmetric
under exchange of two coordinates,
pair permutations that differ by this
exchange; contributions cancel,
giving $\boxed{\Aop\psi=0}$.
The one-electron spin Hilbert space
has only dimension two.
A completely antisymmetric spin state
of $Z$ spin-half electrons would lie
in $\bigwedge^Z\mathbb C^2$, which
vanishes whenever $Z>2$.
Thus $\boxed{Z>2:\text{ no nonzero
completely antisymmetric pure-spin factor}}$.
For atomic fermions the full position--spin
state is antisymmetric, but it typically
cannot be separated into pure spatial and
pure spin factors.

\pnum{5.13}{Autoionization energy and the hydrogenic helium-ion spectrum}
In the independent-electron approximation
of helium ($Z=2$), a hydrogenic electron
has $E_n=-54.4/n^2$ eV.
Two electrons both initially occupying
$n=2$ contribute
$E_i=2(-54.4/4)=-27.2$ eV.
If one falls to $1s$ while the other
is ejected, the final helium ion has
$E_{\mathrm{He}^+}=-54.4$ eV plus the
free electron's kinetic energy $K$.
Energy conservation
$-27.2=-54.4+K$ gives
\begin{equation}
\boxed{K=27.2\ {\rm eV}}.
\end{equation}
The helium ion $\mathrm{He}^+$ has one
electron and therefore the Coulomb
spectrum $E_n=-4(13.6\ {\rm eV})/n^2$.
Its radiative transitions obey
\begin{equation}
\boxed{\frac1\lambda=
4R_H\left(\frac1{n_f^2}-\frac1{n_i^2}\right),
\qquad n_i>n_f}.
\end{equation}
Finite nuclear-mass corrections change
the Rydberg constant slightly.
\checkit{The autoionization estimate ignores
electron--electron repulsion and requires
a physically allowed initial configuration;
it is a textbook idealization.}

\pnum{5.14}{Hypothetical exchange statistics and helium multiplets}
Physical electrons are identical fermions:
spatially symmetric orbitals pair with the
antisymmetric spin singlet, and spatially
antisymmetric orbitals pair with the symmetric
spin triplet. In ground-state helium both
electrons occupy $1s$, hence the spatial part
is symmetric and only a spin singlet is allowed.
If the electrons were instead hypothetical
identical \emph{bosons} with spin $1/2$,
the \emph{total} state would be symmetric.
The two-$1s$ ground configuration would
therefore have the symmetric spin triplet,
with $\boxed{3}$ spin projections.
For excited distinct orbitals $a,b$,
the symmetric spatial sector would pair with
the triplet, and the antisymmetric spatial
sector with the singlet. The former tends
to have larger direct Coulomb repulsion
than the latter; the relative para/ortho
pattern would be interchanged.
If the particles were \emph{distinguishable}
but retained the same masses and charges,
there would be no mandatory spin--space
symmetry pairing. Spatial symmetric and
antisymmetric eigenfunctions could each
occur with all four independent spin
configurations (one singlet and three triplets).
In the spin-independent Hamiltonian,
each spatial energy would have a fourfold
spin degeneracy, aside from any independent
spatial degeneracies. These hypotheticals
alter exchange rules, not the Coulomb operator.

\pnum{5.15}{Helium ground-state Coulomb repulsion integral}
In the uncorrelated $1s^2$ trial state
for charge $Z=2$,
\begin{equation}
\psi_0(\mathbf r_1,\mathbf r_2)
=\frac{Z^3}{\pi a^3}
\exp[-Z(r_1+r_2)/a],\qquad Z=2,
\end{equation}
where $a$ is the ordinary Bohr radius.
The normalized one-particle density is
$\rho(r)=Z^3\exp(-2Zr/a)/(\pi a^3)$.
The angular average of the Coulomb kernel
for fixed $r_1,r_2$ satisfies
\begin{equation}
\frac1{4\pi}\int
\frac{\dd\Omega_2}{|\mathbf r_1-\mathbf r_2|}
=\frac1{\max(r_1,r_2)};
\end{equation}
it follows by integrating
$[r_1^2+r_2^2-2r_1r_2u]^{-1/2}$
over $u\in[-1,1]$.
Set $b=2Z/a$.
The expectation is then
\begin{align}
I&=\int\dd^3r_1\dd^3r_2
\frac{\rho(r_1)\rho(r_2)}{|\mathbf r_1-\mathbf r_2|}\\
&=2b^6\int_0^\infty\dd r_1\,
r_1\ee^{-br_1}
\int_0^{r_1}\dd r_2\,r_2^2\ee^{-br_2}\\
&=\frac{5b}{16}=\frac{5Z}{8a}.
\end{align}
Thus for helium $\boxed{I=5/(4a)}$.
The interaction energy is
\begin{equation}
\Delta E=\frac{e^2}{4\pi\epsilon_0}I
=\frac54\frac{e^2}{4\pi\epsilon_0a}
\simeq\boxed{34.0\ {\rm eV}}.
\end{equation}
Adding to the independent-electron value
$E_0\simeq-108.8$ eV gives
$\boxed{E_{\rm trial}\simeq-74.8\ {\rm eV}}$,
compared with experimental
$E_{\rm He}\simeq-78.98$ eV.
This is a first-order / variational
estimate with an unoptimized $1s^2$ trial
state, not an exact helium solution.

\pnum{5.16}{Lithium: determinant ground state and degeneracy}
Neglect electron--electron interactions
for nuclear charge $Z=3$. The lowest
hydrogenic orbital has energy
$E_{1s}=-13.6Z^2=-122.4$ eV.
Two electrons occupy the distinct spin
orbitals $1s\uparrow$, $1s\downarrow$.
The third occupies $2s$ with
$E_{2s}=-122.4/4=-30.6$ eV.
Therefore
\begin{equation}
\boxed{E_{\rm Li}\simeq
2E_{1s}+E_{2s}=-275.4\ {\rm eV}}.
\end{equation}
For each $\sigma\in\{\uparrow,\downarrow\}$
the normalized wavefunction is
\begin{equation}
\boxed{
\Psi_{\sigma}(1,2,3)=
\frac1{\sqrt{3!}}\det
\begin{pmatrix}
1s\uparrow(1)&1s\downarrow(1)&2s\sigma(1)\\
1s\uparrow(2)&1s\downarrow(2)&2s\sigma(2)\\
1s\uparrow(3)&1s\downarrow(3)&2s\sigma(3)
\end{pmatrix}}.
\end{equation}
These are eigenstates of the sum of
hydrogenic one-body Hamiltonians and
are fully antisymmetric. The closed
$1s^2$ shell has zero total spin;
the $2s$ valence electron provides
a spin doublet:
$\boxed{g_{\rm ground}=2}$.
